\subsection*{Why the exchange exists}
The weighted criterion also explains the failure of the conjecture without a
choice of particular digit counts. More generally let $A_0$ be any
three-term-arithmetic-progression-free subset of the nonnegative integers
with $|A_0\cap[0,x]|\ll (1+x)^\alpha$ for some $0\le\alpha<1$.
Its deficit rule has limiting ball density
$(1-q^{-1})\sum_{d\in A_0}q^{-d}$. For every $q$ one can construct a
polynomial-ratio progression-free set of strictly greater upper density.

Indeed, for $X\ge1$, Abel summation gives
$\sum_{d\in A_0}e^{-d/X}\ll_\alpha X^\alpha$, with a constant also
depending on the counting bound. Choose an integer $h\ge2$ so large that
$h>(2h-1)^\alpha$. In dimension $n$ read the digits $0,\ldots,h-1$
in base $2h-1$ and set $X=(2h-1)^n$. One squared-norm level contains at
least $h^n/(1+n(h-1)^2)$ vectors. Its integer encodings form a
three-term-free set $B$ in $[0,X/2]$ by the same carry-free midpoint
argument as before. Consequently
\[
 \sum_{d\in B}e^{-d/X}\ge
 \frac{e^{-1/2}h^n}{1+n(h-1)^2}>\sum_{d\in A_0}e^{-d/X}
\]
for all sufficiently large $n$. Use $z=e^{-1/X}$ in the positive-coefficient
argument for the weighted exchange, with $A_0$ in place of $A$. The proof
of confinement only uses the finiteness of $B$ and progression avoidance
inside and outside the modified classes. The same separated degree blocks
give the infinite set.

For the ternary $A$ in the conjecture one can take $\alpha=\log_3 2$,
because there are $2^k$ permitted integers below $3^k$. Already $h=4$
satisfies $4>7^{\log_3 2}$. Thus the transferable ingredient is a growth
gap between two explicit digit constructions. The numerical certificates
give finite versions of this argument.
